% !TEX program = lualatex
\documentclass{article}
\usepackage{geometry}\geometry{margin=23mm}
\usepackage[bidi=basic]{babel}\babelprovide[import,main]{arabic}
\babelfont{rm}[ItalicFont={Times New Roman},ItalicFeatures={FakeSlant=.13}]{Times New Roman}
\begin{document}
\title{لوح القياس}\author{}\date{١٤ مايو ٢٠٢٦}\maketitle
\begin{abstract}تدل علامتان على موضع اللوح الأول.\end{abstract}
\tableofcontents
\section{العلامات}\begin{quote}«اقرأ العلامة السفلى أولًا».\end{quote}
\begin{figure}[ht]\centering\fbox{\rule{0pt}{8mm}\rule{33mm}{0pt}}\caption{موضع اللوح}\end{figure}
\section{الارتفاع}
\begin{table}[ht]\centering\begin{tabular}{cr}العلامة&الارتفاع (mm)\\\hline 1&12\\2&18\end{tabular}\caption{ارتفاع العلامات}\end{table}
\begin{thebibliography}{1}\bibitem{record}\emph{دفتر القياس}, 2025, \chaptername~2.\end{thebibliography}
\end{document}
